% =============================================================================
% INTRODUCTION TO VOLUME 5: LAYERS, HILBERT-SPACE BRIDGE AND RECENT CLARIFICATIONS
% (Vol. 5 of 5; continuation of Vol. 4)
% =============================================================================

\section*{Fifth Volume of the Document Collection}

This fifth volume concludes the collection of individual documents on T0
theory / FFGFT. It gathers the works that carry the methodological
consolidation of the theory and, to a particular degree, mark its present
state.

\subsection*{Relation to Volume 4}

Volumes 4 and 5 together replace the single extension volume originally
planned, which would have exceeded the 550-page Kindle limit. The transition
from Volume~4 to Volume~5 is not arbitrary but follows a content boundary:
up to Doc.~184 (p-bit), the theory is still developed essentially on
Layer~1 (compact $T^4$, correction-free description in natural units); from
Doc.~185 (Embedding Cost) onwards, the Layer-1 / Layer-2 distinction is
used as an explicit descriptive grid and carried through to the epistemic
self-positioning of Doc.~262.

\subsection*{Character of this Volume}

Volume 5 is denser than the earlier volumes because its documents form the
methodological-foundational scaffolding of the most recent development phase:
\begin{itemize}
\item The Hilbert-space bridge (Doc.~230--232), which formally relates FFGFT
      to quantum mechanics and articulates operational equivalence with
      interpretive divergence.
\item The layers and scale-ladder documents (Doc.~241--253), which
      systematically work through the description in Layer~1 (static,
      compact) and Layer~2 (dynamic, SI-projected).
\item The most recent clarifications (Doc.~255--262), culminating in the
      epistemic self-positioning of Doc.~262 (acceptance without
      visualisation and inside view of decompactified time).
\end{itemize}

\subsection*{Five-Part Organisation}

The 37 documents are arranged in numerical order and divided into five
thematic parts:

\begin{itemize}
\item \textbf{Part I --- Late Geometry and Layer Preparation}: embedding
      cost, FFGFT photonics with $\mathrm{LiNbO}_3$, geometric foundations,
      torus derivation, corrections register, Gemini dialogue, algebraic
      composition limits, non-closure theorem.
\item \textbf{Part II --- FFGFT Field Theory and Operators}: complete
      field-theory synthesis, recursion operator $\mathcal{R}$, fractal
      boundary condition, narrative, triangle-matrix reduction, consciousness
      bridges, corrections and windings.
\item \textbf{Part III --- Hilbert-Space Bridge}: Hilbert-space translation,
      Hilbert-space extension, QG as a subset.
\item \textbf{Part IV --- Layers and Information Formalism}: two layers,
      scale ladder, information formalism, UMF analysis, RA diagnostic,
      RSG comparison, information as state and process, epistemology,
      falsifiability, information and SL, comparisons with Vopson and
      Phillips, xi number-range analysis.
\item \textbf{Part V --- Recent Clarifications and Epistemic
      Self-Positioning}: information from kinetics, information unit and
      scale, Koide factor 2/3, Koide cross-terms, static natural units and
      SI projection, and acceptance without visualisation and inside view of
      decompactified time (Doc.~262 as the methodologically foundational
      closing chapter).
\end{itemize}

\subsection*{Status}

All documents in this volume carry the May 2026 status.
